\documentclass[conference]{assets/latex/IEEEtran}

\usepackage{amsmath,amssymb,amsfonts}
\usepackage{algorithmic}
\usepackage{graphicx}
\usepackage{textcomp}
\usepackage{xcolor}
\usepackage[backend=biber,style=numeric]{biblatex}
\addbibresource{assets/latex/report.bib}

\graphicspath{{assets/latex/}}

\begin{document}

\title{CS6524 Computer Vision Lab. 3\\ Structure from Motion}

\author{\IEEEauthorblockN{1\textsuperscript{st} Frederick Sheppard}
\IEEEauthorblockA{\textit{23361433} \\
\textit{4th Year ISE}}
\and
\IEEEauthorblockN{2\textsuperscript{nd} Oluwajomiloju Olajitan}
\IEEEauthorblockA{\textit{23373326} \\
\textit{4th Year ISE}}
\and
\IEEEauthorblockN{3\textsuperscript{rd} Tiernan Scully}
\IEEEauthorblockA{\textit{23365528} \\
\textit{4th Year ISE}}
}

\maketitle

\begin{abstract}
TODO
\end{abstract}




\section{Methodology}\label{sec:methodology}

\subsection{Input and Calibration}
Video (a) moves the camera around a plug adaptor and video (b) turns the
plug in front of a still camera. Frames are rotated to landscape and
resized to $640\times360$. We reuse the Lab 2 calibration of the same
phone model, scaling $K$ to the new resolution. The original distortion
fit folded near the image corners, so we refit it with one radial
coefficient.

\subsection{Feature Tracking}
SIFT~\cite{lowe2004sift} had a higher RANSAC inlier fraction than ORB on
both videos (0.99 against 0.90 on (a)) and gave half the reprojection
error in a full run, so we use it throughout. We keep 2000 keypoints per
frame, match consecutive frames with Lowe's ratio test at 0.75, and pass
track IDs along matched points. In (b) most matches were on the static
background, so we restrict detection to a region around the plug and
drop matches that move less than 0.4 pixels.

\subsection{Pose and Triangulation}
For each consecutive pair we estimate $E$ with the five point
algorithm~\cite{nister2004five} in RANSAC and recover $R$ and $t$ with a
cheirality check. Points are triangulated with the first camera at the
origin and kept if they are in front of both cameras, reproject within
2 pixels, and have at least 0.5 degrees between their rays.

\subsection{Extending to N Frames}
Consecutive frames barely move, so we work on keyframes (every 10th frame
for (a), every 3rd for (b)). New keyframes are placed against the map
with PnP RANSAC and add new triangulated points. Every 5 keyframes we run
bundle adjustment~\cite{triggs2000bundle} over the last 5, and once over
all keyframes at the end, minimising
\begin{equation}
\sum_{i,j} \rho\left(\left\lVert x_{ij} - \pi(K, R_i, t_i, X_j)\right\rVert^2\right)
\end{equation}
with a soft L1 loss $\rho$. Points seen in fewer than 3 keyframes or with
median error above 1.5 pixels are removed.

\subsection{Dense Reconstruction}
We pick up to 20 keyframe pairs, rectify them using their poses, and
compute disparity with semi-global matching~\cite{hirschmuller2008sgm}.
Points are kept if they land inside the segmented plug in most views and
at least two pairs agree on their depth within 2\%. Points in the same
voxel are averaged and saved as a PLY.

\subsection{Sensitivity Analysis}
The full pipeline is rerun with $f_x, f_y$ scaled by 0.8, 1.0 and 1.2,
with and without undistortion, using a fixed random seed so only the
camera model changes.

\section{Results}\label{sec:results}
Table~\ref{tab:results} shows the twelve runs. Video (a) registered all 147
keyframes in every run and kept around 16,000 sparse points. Its median
reprojection error was lowest at a focal factor of 1.0 (0.20 pixels
undistorted, 0.18 distorted) and rose at both 0.8 and 1.2, which suggests the
reused $K$ is close to the true focal length. Removing undistortion did not
increase the error, since $k_1$ is small and bundle adjustment can absorb some
of the lens error into the poses and points. These errors are measured on the
points the pipeline kept, so they show consistency, not 3D accuracy.

Video (b) is much weaker. At most 63 of 79 keyframes were registered and only
127 to 331 sparse points survived, mostly on the pins and cable socket. At a
factor of 1.0 the map only covers frames 39 to 210, so the full rotation is not
recovered. The dense point count also does not follow quality: b\_12d has the
most dense points but the fewest registered views, and a\_12u kept only 854
points while a\_12d kept 12,101 from almost the same sparse map.

\begin{table}[htbp]
\caption{Sensitivity runs. Views are registered keyframes, error is median reprojection error in pixels.}
\label{tab:results}
\centering
\begin{tabular}{l r r r r}
\hline
Run & Views & Sparse & Dense & Error \\
\hline
a\_08u & 147 & 15,440 & 3,666 & 0.26 \\
a\_08d & 147 & 16,199 & 6,904 & 0.22 \\
a\_10u & 147 & 16,946 & 14,570 & 0.20 \\
a\_10d & 147 & 16,174 & 15,633 & 0.18 \\
a\_12u & 147 & 16,750 & 854 & 0.22 \\
a\_12d & 147 & 15,991 & 12,101 & 0.20 \\
b\_08u & 61 & 312 & 27,030 & 0.26 \\
b\_08d & 50 & 235 & 61,364 & 0.25 \\
b\_10u & 57 & 258 & 78,455 & 0.25 \\
b\_10d & 63 & 331 & 67,837 & 0.26 \\
b\_12u & 54 & 209 & 20,358 & 0.29 \\
b\_12d & 27 & 127 & 99,215 & 0.32 \\
\hline
\end{tabular}
\end{table}

\section{Discussion \& Failure Analysis}\label{sec:discussion}



\section{Conclusion}
We built a full structure from motion pipeline that tracks SIFT features,
estimates poses from the essential matrix, corrects drift with bundle
adjustment and produces dense point clouds for both videos under six camera
settings. The camera orbit worked well, registering every keyframe with sub
pixel error, while the spinning object was limited by its white surface and
static background. The main lessons were that consecutive frames give too
small a baseline for reliable poses, and that a low reprojection error does not
guarantee a correct shape. With more time, calibrating the actual phone and
segmenting the object before feature detection would be the most useful
improvements.

\printbibliography

\end{document}
